%% ============================================================
%% Draft section: Forgetting in LiB — ordinal re-ranking (active) +
%% probation/life (passive), matching trainer.rs/trie.rs
%% Companion to paper.tex. Numbers marked [EN] are measured on English
%% (Wikipedia, 50k cap, 10k document-epochs, doc_size=50); \TODO marks
%% values awaiting the multilingual re-run.
%%
%% Intended placement: expands §3.2 (which currently describes admission
%% only) and adds a short results subsection. Two blocks below:
%%   (A) algorithm description — replaces the admission-only §3.2 text;
%%   (B) results — the forgetting characterisation.
%% ============================================================

\newcommand{\TODO}[1]{\textcolor{todocolor}{[\textbf{TODO}: #1]}}

%% ------------------------------------------------------------
%% (A) Algorithm: memorising AND forgetting
%% ------------------------------------------------------------

\subsection{Memorising and Forgetting}
\label{sec:forgetting}

LiB follows the Principle of Least Effort by keeping the lexicon $L$ small: it
must \emph{forget} units as readily as it memorises them, or the lexicon grows
without bound and $\mathrm{DL}_\mathrm{lex}$ is never controlled.
\citet{yang2020less} accordingly build forgetting into the model as a
first-class mechanism with two routes — an \emph{active} route that demotes
units judged unhelpful, and a \emph{passive} route that lets unused units decay
— mirroring the active and passive forgetting observed in human memory.
We make both routes explicit here, as together they are central to how LiB
adapts its vocabulary size to a corpus.

The lexicon $L$ is a priority-ordered list: each unit $u$ has an ordinal
position $\Theta(u)$, with $\Theta=0$ the head. Segmentation itself never
consults the ranking — it is greedy longest-match with a one-chunk lookahead —
so $\Theta$ functions purely as a memory: it encodes how strongly a unit is
currently held, and it decides who is at risk of being forgotten.
Training proceeds over \emph{documents}: each epoch draws one document of
$d$ sentences ($d=\texttt{doc\_size}$). Within a document:

\begin{description}
  \item[Read \& memorise.] Each sentence is segmented with the current $L$;
    with probability $\alpha$ (\texttt{memory\_in}) the concatenation of two
    adjacent chunks is proposed as a candidate unit (an unknown character
    sequence of length ${\le}2$ may itself be a candidate). A candidate
    sighted \emph{twice within the same document} is admitted, provided
    $|L|$ is below the cap $|V|$. New units enter at the tail — the lowest
    rank — and must earn their way up.
  \item[Evaluate (active forgetting).] At each position the segmenter also
    considers the second-longest match $c_2$ against the longest $c_1$; both
    branches are rolled out greedily until they reconverge, and $c_1$ is
    judged \emph{bad} if its branch needs more unknown symbols, or (at equal
    unknowns) more chunks, or (at equal both) a higher ordinal sum
    \citep[\S 2.3.1]{yang2020less}. A good $c_1$ is promoted headward,
    $\Theta\leftarrow\Theta-(\lfloor\Theta\Delta\rfloor+1)$ with
    $\Delta=\texttt{update\_rate}$; a bad $c_1$ is demoted tailward by the
    same step — and a unit demoted \emph{past} the tail ($\Theta\ge|L|$) is
    deleted outright — while the longest match of its shortened form is
    promoted in its place \citep[\S 2.3.2]{yang2020less}. Promotion is thus
    multiplicative: units that keep winning comparisons converge
    exponentially fast on the head, where they are safe.
  \item[Probation \& expiry (passive forgetting).] Once per document, every
    unit on the probation watch has its remaining life decremented, and units
    whose life expires are forgotten. Then the current tail — the bottom
    fraction $\omega$ (\texttt{memory\_out}) of $L$ — is placed on watch with
    life $\lambda$ (\texttt{life}); seed units start on watch. Probation is
    cancelled the moment a unit is used in a segmentation without being judged
    bad, so
    only units that go $\lambda$ consecutive documents without proving
    themselves are lost.
\end{description}

Atomic units (single characters, byte-fallback tokens) and special tokens are
exempt from both deletion routes, guaranteeing lossless reconstruction.
The two routes play complementary roles: active forgetting is
\emph{comparative} — it punishes units that lose against an available
alternative, immediately, wherever they rank — while passive forgetting is
\emph{absolute} — it removes units that are simply never exercised, after a
grace period of $\lambda$ documents. A unit that stops being useful is first
demoted out of the safe head region, then drifts into the watched tail, and
finally expires.

Crucially, the vocabulary size is \emph{not} imposed by the cap $|V|$: it is an
emergent equilibrium between admission and the two forgetting routes, governed
by $\omega$.
This is the mechanism behind the language-sensitive saturation reported in
Section~\ref{sec:multilingual}: rather than an accident of one language reaching
a wall, saturation is the general behaviour of the model once forgetting is
strong enough to balance admission.

%% ------------------------------------------------------------
%% (B) Results: forgetting controls vocabulary self-regulation
%% ------------------------------------------------------------

\subsection{Forgetting Controls Vocabulary Self-Regulation}
\label{sec:forgetting-results}

The forgetting ratio $\omega$ is the knob that sets where the
admission--forgetting equilibrium falls.
Table~\ref{tab:omega} sweeps $\omega$ on English at a 50k cap and reports the
resulting (emergent) vocabulary size, the supra-word ratio, and the fraction of
supra-word types that never occur when re-encoding the training corpus — a
direct measure of \emph{dead weight}, units the lexicon stores but never uses.

\begin{table}[t]
\centering
\small
\begin{tabular}{rrrr}
\toprule
$\omega$ & \textbf{$|L|$} & \textbf{supra \%} & \textbf{never-used \%} \\
\midrule
$1\times10^{-4}$ & 50{,}000 & 56.5 & 5.9 \\
$1\times10^{-3}$ & 49{,}955 & 53.1 & 3.2 \\
$5\times10^{-3}$ & 49{,}755 & 50.8 & 1.9 \\
$2\times10^{-2}$ & 21{,}002 & 37.1 & 0.5 \\
\bottomrule
\end{tabular}
\caption{Effect of the forgetting ratio $\omega$ on English (50k cap; sweep
run at 10k document-epochs with $d{=}50$ — an earlier configuration than the
headline runs, whose qualitative behaviour it characterises). As forgetting
strengthens, dead weight (never-used supra-word types) falls monotonically,
and beyond a threshold the vocabulary stops filling the cap and
self-regulates to a far smaller emergent size.}
\label{tab:omega}
\end{table}

Two effects are visible.
First, dead weight falls monotonically with $\omega$: the never-used supra-word
fraction drops from $5.9\%$ to $0.5\%$, because stronger forgetting continually
evicts the least-used units and fills the freed budget with units that earn
their place. Even in the regime where the lexicon remains full, its
\emph{composition} improves.
Second, there is a sharp transition between $\omega=5\times10^{-3}$ and
$\omega=2\times10^{-2}$: below it, admission outpaces forgetting and the lexicon
pins at the cap; above it, forgetting wins and the lexicon self-regulates far
below the cap (from 50k to 21k for English). The near-discontinuity marks the
point where the per-document eviction rate overtakes the admission rate.

This reframes the vocabulary size as a \emph{result} rather than a
hyperparameter. Standard tokenizers fill whatever $|V|$ they are given; LiB, with
forgetting active, discovers the vocabulary size at which no further unit repays
its storage cost, and $\omega$ tunes how aggressively it does so. The supra-word
ratio likewise becomes a controllable, emergent quantity rather than a fixed
property, which we relate to typology in Section~\ref{sec:multilingual}.

\paragraph{Reproducibility note.}
The two forgetting routes are the load-bearing part of the LiB training loop.
All LiB results in this paper use the implementation in which active
re-ranking and passive probation are both operating, as described above, with
the per-language configurations of Section~\ref{sec:setup}; we release this
implementation with the paper.
